\DocumentMetadata{
  pdfversion=2.0,
  pdfstandard=ua-2,
  lang=en-US,
  tagging=on,
  tagging-setup={math/setup=mathml-SE,tabsorder=structure}
}
\documentclass[11pt,letterpaper]{article}
\usepackage[margin=0.68in,includefoot]{geometry}
\usepackage{newcomputermodern}
\usepackage{amsmath,amscd,mathtools}
\usepackage{tikz,adjustbox}
\providecommand{\square}{\mdlgwhtsquare}
\usepackage[hidelinks]{hyperref}
\hypersetup{
  pdftitle={Algebra First-Year Exam, August 2020, Part 2},
  pdfauthor={University of Florida Department of Mathematics},
  pdfsubject={Graduate mathematics examination},
  pdfdisplaydoctitle=true,
  bookmarksopen=true
}
\setcounter{secnumdepth}{0}
\setlength{\parindent}{0pt}
\setlength{\parskip}{0.28em}
\setlength{\emergencystretch}{3em}
\setlength{\leftmargini}{2.15em}
\setlength{\leftmarginii}{2.65em}
\setlength{\leftmarginiii}{3em}
\renewcommand{\labelenumi}{\textbf{\theenumi.}}
\pagestyle{empty}
\raggedbottom
\allowdisplaybreaks
\newenvironment{examproblems}[1]{%
  \begin{enumerate}%
  \setcounter{enumi}{#1}%
  \setlength{\topsep}{0.35em}%
  \setlength{\partopsep}{0pt}%
  \setlength{\itemsep}{0.48em}%
  \setlength{\parsep}{0.18em}%
}{\end{enumerate}}
\newenvironment{examsubparts}{%
  \begin{enumerate}%
  \setlength{\topsep}{0.18em}%
  \setlength{\partopsep}{0pt}%
  \setlength{\itemsep}{0.16em}%
  \setlength{\parsep}{0.1em}%
}{\end{enumerate}}
\newenvironment{examinstructions}{%
  \par\begingroup\itshape
}{%
  \par\endgroup\vspace{0.18em}
}
\makeatletter
\renewcommand\section{\@startsection{section}{1}{0pt}%
  {-1.25ex plus -.25ex minus -.1ex}%
  {0.55ex plus .1ex}%
  {\normalfont\large\bfseries}}
\renewcommand{\@maketitle}{%
  \newpage\null\vskip -1em%
  {\footnotesize\bfseries
    \href{https://www.ufl.edu/}{University of Florida}, \href{https://math.ufl.edu/}{Department of Mathematics}\par}%
  \vskip -0.5em%
  \tagstructbegin{tag=Title}%
    {\LARGE\bfseries\href{https://gma.math.ufl.edu/exams/algebra-first-year/}{Algebra First-Year Exam}, August 2020, Part 2\par}%
  \tagstructend%
  \vskip 0.25em%
  \tagmcbegin{artifact}%
    \hrule height 1pt%
  \tagmcend%
  \vskip 1em%
}
\makeatother
\title{Algebra First-Year Exam, August 2020, Part 2}
\author{}
\date{}
\begin{document}
\maketitle
\thispagestyle{empty}
\begin{examinstructions}
Answer four problems. Write your answers clearly in complete English sentences. You may quote results (within reason) as long as you state them clearly.
\end{examinstructions}
\begin{examproblems}{0}
\item
Suppose~\(R\) is a commutative ring with identity and~\(F\) is a free~\(R\)-module of rank~\(n\). Show that for any~\(R\)-module~\(M\) the~\(R\)-module \(\operatorname{Hom}_R(F,M)\) is isomorphic to the direct sum of~\(n\) copies of~\(M\).
\item
Which of the following polynomials are irreducible? Explain, and factor the ones that are not irreducible.
\begin{examsubparts}
\item[\textnormal{(a)}]
\(X^3+X+1\) in \(\mathbb{Q}[X]\);
\item[\textnormal{(b)}]
\(X^4+X^2+1\) in \(\mathbb{Q}[X]\);
\item[\textnormal{(c)}]
\(Y^2-X^3+X^2\) in \(\mathbb{C}[X,Y]\).
\end{examsubparts}
\item
Let~\(R\) be an integral domain.
\begin{examsubparts}
\item[\textnormal{(a)}]
Define what it means for an element of~\(R\) to be irreducible.
\item[\textnormal{(b)}]
Define what it means for an element of~\(R\) to be prime.
\item[\textnormal{(c)}]
Show that a prime element is irreducible.
\end{examsubparts}
\item
Suppose~\(F\) is a field, \(V\) is an~\(F\)-vector space and \(f:V\to V\) is a linear map.
\begin{examsubparts}
\item[\textnormal{(a)}]
Show that if~\(V\) has finite dimension then~\(f\) is injective if and only if it is surjective.
\item[\textnormal{(b)}]
Give examples with~\(V\) of infinite dimension, in which~\(f\) is surjective but not injective.
\item[\textnormal{(c)}]
Give an example in which~\(f\) is injective but not surjective. Justify both examples; you may use any results on dimension if they are stated clearly.
\end{examsubparts}
\item
A minimal prime ideal of a ring~\(R\) is a prime ideal that does not properly contain any other prime ideal. If~\(F\) is a field, find all minimal prime ideals in the ring
\[
R=F[x_1,x_2,x_3,x_4,x_5,x_6]/(x_1x_2,x_3x_4,x_5x_6).
\]
\item
Suppose~\(K\) is a field and \(a\in K\) is such that \(X^3-a\) and \(X^2+X+1\) are irreducible in \(K[X]\). If \(L=K[X]/(X^3-a)\), show that~\(L\) contains exactly one root of the polynomial \(X^3-a\).
\item
This problem has two parts.
\begin{examsubparts}
\item[\textnormal{(a)}]
Find representatives in rational canonical form for all conjugacy classes in \(GL_6(\mathbb{Q})\) with characteristic polynomial \((X^3-1)^2\).
\item[\textnormal{(b)}]
For each representative, give the Jordan normal form in \(GL_6(\mathbb{C})\).
\end{examsubparts}
\end{examproblems}
\end{document}
